\section{A bound for the row determinant}\label{sec:signed}

This section proves Proposition~\ref{prop:signed-goal}, which completes the proof of Theorem~\ref{thm:real}.
In the case with doubled interaction OpenAI's integrand contains a determinant $\det\mathsf L$ of size $n$
that the energy method bounds by $e^{O(n\log n)}$. We replace this determinant by a simpler one, with a loss of only
$e^{O(N)}$, and we bound the simpler one by a Gram determinant. The argument has three steps.

The first step is a comparison. Let $F(w)$ be the Pick matrix of the nodes with weight $w$ on the positive
nodes. Then $\det\mathsf L$ is $\det F(-5)$ up to a known factor, and a positivity argument bounds it by
$499^n\det F(0)$ as soon as $F(-\frac1{100})\ge0$ (\S\ref{ssec:rowdet}). The second step shows that
$F(-\frac1{100})\ge0$ whenever a phase inequality holds for the nodes (\S\ref{ssec:window},
\S\ref{ssec:windowproof}), and that this inequality holds on every configuration for which the energy bound
is not already sufficient (\S\ref{ssec:alternative}). The third step bounds $\det F(0)$, which is a Vandermonde
determinant of a Laurent system. It splits into a Vandermonde determinant of the positive nodes and a
determinant of folded functions at the negative nodes. Fischer's inequality and Andr\'eief's identity turn
the latter into two log-gas integrals (\S\ref{ssec:gram}). Reference measures then bound these integrals,
and a numerical inequality closes the proof (\S\ref{ssec:budget}).

\subsection{The row determinant and the comparison}\label{ssec:rowdet}
In the case with doubled interaction OpenAI's integrand contains the determinant of the row functions at
points $z_1,\ldots,z_n\in(-1,1)\setminus\{0\}$, which we call the nodes \cite[Prop.~5.2]{OpenAI2026}. They are
the free points $x_i$ and the prescribed points $y_j$ together. We write $\det\varphi(z)$ for this determinant. With $\mathsf e=n-1-2g$ it factors as
\begin{equation}\label{eq:detphi}
 \det\varphi(z)=C_{n,g}\,V(z)\prod_a|z_a|^{-(n-1-g)}\det\mathsf L(z),
\end{equation}
where $C_{n,g}$ depends only on $n$, $g$ and the number of nodes of each kind, and $V(z)=\prod_{a<b}(z_b-z_a)$. To describe $\mathsf L$ let
\[
 B(z)=\prod_a\frac{z-z_a}{1-z_az},\qquad \cK_{ab}=\frac1{1-z_az_b},\qquad D_a=B'(z_a)z_a^{\mathsf e},
\]
so that $B$ is the Blaschke product of the nodes and $\cK$ is the Pick matrix of the Szeg\H{o} kernels. Let
$E_-$ and $E_+$ be the diagonal blocks of $\operatorname{diag}(D_a)$ at the negative and the positive nodes.
For $w\in\R$ let
\[
 F(w)=\cK+\operatorname{diag}(E_-,wE_+).
\]
The computation of \cite[Prop.~5.2]{OpenAI2026} gives, for every real vector $(w_a)$, the identity
\begin{equation}\label{eq:laurent}
 \det\bigl(z_a^{g-r}+w_az_a^{r-g}\bigr)_{a,\,0\le r<n}=c(z)\det\bigl(\cK+\operatorname{diag}(w_aD_a)\bigr),
\end{equation}
where $c(z)$ does not depend on $w$. If all nodes are free, the rows of OpenAI's determinant are
$\frac12(z^{g-r}+z^{r-g})$ at a negative node and $-\frac14z^{g-r}+\frac54z^{r-g}$ at a positive node. So
$\det\mathsf L=\det F(-5)/\det\cK$. The value $-5=1-6$ comes from the factor $\frac32$ in the row functions.
OpenAI's argument uses only $|\det\mathsf L|\le e^{O(n\log n)}$, which is negligible at the scale $e^{N^2}$. Numerical
evaluation suggests that $\log|\det\mathsf L|$ is of order $-N^2$ near the extremal configuration
(Remark~\ref{rem:L}).

A prescribed point has the row $\frac12(z^{g-r}+z^{r-g})$ at either sign. At a positive prescribed point this row
is $-2$ times the row of a free positive node plus $3$ times the row $(z^{r-g})_r$. Let $\Xi$ be the matrix of
all rows at the nodes. Expanding its determinant in these rows gives
\begin{equation}\label{eq:expandS}
 |\det\Xi|\le\sum_{S}2^{m_+-|S|}3^{|S|}|\det\Xi_S|,
\end{equation}
where $m_+$ is the number of positive prescribed points and $S$ runs over the subsets of these points. In
$\Xi_S$ a node of $S$ has the row $(z^{r-g})_r$, a negative node has the row $\frac12(z^{g-r}+z^{r-g})$,
and every other positive node has the row of a free positive node. The total weight in \eqref{eq:expandS} is at
most $5^m$. A node of $S$ is what Section~\ref{sec:energy} calls a prescribed point of the second kind. Let
$\Xi^0_S$ be the same matrix with the rows of the positive nodes outside $S$ replaced by $(z^{g-r})_r$.

Both $\Xi_S$ and $\Xi^0_S$ are obtained from \eqref{eq:laurent}. The left side of \eqref{eq:laurent} is
affine in each $w_a$. Taking the coefficient of $\prod_{a\in S}w_a$ gives
\begin{equation}\label{eq:minors}
 \frac{|\det\Xi_S|}{|\det\Xi^0_S|}=4^{-|P\setminus S|}\frac{|\det F(-5)_{S^c}|}{\det F(0)_{S^c}},
\end{equation}
where $P$ is the set of positive nodes and the index $S^c$ denotes the principal submatrix on the complement of
$S$. The case $w=1$ is a Pick matrix.

\begin{lemma}\label{lem:pick}
Let $H=\operatorname{diag}(B'(z_a))$ and let $X$ be a real diagonal matrix. If $\cK-X\cK X\ge0$, then $\cK+HX\ge0$. In
particular $F(1)\ge0$ for every configuration.
\end{lemma}

\begin{proof}
First we show that $\cK H^{-1}\cK=H$. For a real vector $c$ put $\phi_c(z)=B(z)\sum_ac_a/(z-z_a)$. It has node values
$\phi_c(z_a)=c_aB'(z_a)$. Partial fractions give $\phi_c=\sum_ad_ak_a$ with $k_a(z)=1/(1-z_az)$, for a real vector $d$, and comparing node values
gives $\cK d=Hc$. On the unit circle $|\phi_c(z)|=|\sum_ac_a/(z-z_a)|=|\sum_ac_ak_a(\bar z)|$. So $\|\phi_c\|_{H^2}^2=c^{\mathsf T}\cK c$, and also
$\|\phi_c\|_{H^2}^2=d^{\mathsf T}\cK d=c^{\mathsf T}H\cK^{-1}Hc$. Hence $H\cK^{-1}H=\cK$, which is the claim.

Let $T=H^{-1}\cK X$ and $A=I+T$. Then $\cK T=HX$, $T^{\mathsf T}\cK=XH$ and $T^{\mathsf T}\cK T=X\cK X$. So
\[
 0\le A^{\mathsf T}\cK A=\cK+2HX+X\cK X .
\]
Adding $\cK-X\cK X\ge0$ gives $2(\cK+HX)\ge0$. For the second statement take $X=\operatorname{diag}(z_a^{\mathsf e})$. The entries of
$\cK-X\cK X$ are $\sum_{j<\mathsf e}(z_az_b)^j$, which is a Gram matrix.
\end{proof}

\begin{lemma}\label{lem:compare}
If $F(-\frac1{100})\ge0$, then $|\det F(-5)_T|\le499^{|T|}\det F(0)_T$ for every set $T$ of nodes. Hence
$|\det\Xi_S|\le499^n|\det\Xi^0_S|$ for every $S$.
\end{lemma}

\begin{proof}
Put $Y=F(0)$ and $Q=\operatorname{diag}(0,E_+)$. By hypothesis and Lemma~\ref{lem:pick}, $Y+Q\ge0$ and
$Y-\frac1{100}Q\ge0$. A convex combination of the two gives $Y\ge0$. Let $Z=F(-5)=Y-5Q$. Then
\[
 499Y-Z=5(Y+Q)+493Y\ge0,\qquad 499Y+Z=500\bigl(Y-\tfrac1{100}Q\bigr)\ge0 .
\]
So $-499Y\le Z\le499Y$, and the same holds for principal submatrices. If $Y_T>0$, all eigenvalues of
$Y_T^{-1/2}Z_TY_T^{-1/2}$ have modulus at most $499$, which is the claim. The singular case follows by replacing
$Y_T$ by $Y_T+tI$ and letting $t\to0$. The second statement follows from \eqref{eq:minors}.
\end{proof}

So everything reduces to two questions: when is $F(-\frac1{100})\ge0$, and how large is $\det\Xi^0_S$?

\subsection{A phase criterion for positivity}\label{ssec:window}
The matrix $F(-c)$ is a Pick matrix with a small negative weight on the positive nodes. The following theorem
gives a criterion for its positivity that depends only on the positions of the positive nodes. It holds for
all configurations and has no relation to our parameters.

\begin{theorem}\label{thm:window}
Let $z_1,\ldots,z_n\in(-1,1)\setminus\{0\}$ be distinct, with $k_+$ positive and $k_-$ negative entries. Let
$t\ge0$ be an integer, let $c=\frac1{100}$, and let $D_a=B'(z_a)z_a^t$. For a positive node $p$ put
$\rho_p=(1-p)/(1+p)$. If
\begin{equation}\label{eq:window}
 t\arctan4-\sum_{p>0}\arctan(4\rho_p)>\frac{\pi}{2}(k_-+1)+\frac{17}4\,\frac{c^2}{1-c^2}\,k_+,
\end{equation}
then $\cK+\operatorname{diag}(E_-,-cE_+)\ge0$.
\end{theorem}

The number $\arctan4$ is the half-angle of the arc of the unit circle around $z=-1$ with end point
$z_0=(4i-1)/(4i+1)$. The proof compares two phases on this arc. We apply the theorem with $t=\mathsf e$.

\subsection{Proof of Theorem~\ref{thm:window}}\label{ssec:windowproof}
Let $f_a(\gamma)=\gamma z_a^t$ at a negative node and $f_a(\gamma)=-cz_a^t$ at a positive node. Let
$W(\gamma)_{ab}=(1-f_a(\gamma)f_b(\gamma))\cK_{ab}$ be the Pick matrix of these data. By Lemma~\ref{lem:pick} it suffices to prove
$W(1)\ge0$. With $X=\operatorname{diag}(z_a^t)$ we have
\[
 W(-c)=(1-c^2)\cK+c^2(\cK-X\cK X),
\]
and $\cK-X\cK X\ge0$ as in Lemma~\ref{lem:pick}. So $W(-c)>0$. If $W(1)$ is not positive semidefinite, then there is
$\gamma\in(-c,1)$ with $W(\gamma)\ge0$ singular. Fix such a $\gamma$ and write $W=W(\gamma)$, $f_a=f_a(\gamma)$. Note that $|f_a|<1$.

Let $\mathcal T$ be a minimal set of nodes such that the principal submatrix $W_{\mathcal T}$ is singular. Since the diagonal entries of $W$ are
positive, $|\mathcal T|\ge2$. By minimality $W_{\mathcal T}$ has a kernel vector $v$ with all entries nonzero, and every proper principal
submatrix is positive definite. From now on only the nodes of $\mathcal T$ are used, and $k_+$ and $k_-$ count these nodes.
Inequality \eqref{eq:window} still holds, because removing a node lowers the right side or raises the left side. If all nodes of $\mathcal T$ have
the same sign, the data are a constant of modulus less than $1$ times $z^t$, and $W_{\mathcal T}>0$ as above. So $k_+\ge1$ and $k_-\ge1$.

We construct a Blaschke product that interpolates the data. Put
\[
 \varphi(\zeta)=\sum_{a\in\mathcal T}\frac{v_a}{1-z_a\zeta},\qquad \psi(\zeta)=\sum_{a\in\mathcal T}\frac{v_af_a}{1-z_a\zeta}.
\]
The relation $Wv=0$ says that $\varphi(z_b)=f_b\psi(z_b)$ for $b\in\mathcal T$. For $|\zeta|<1$ let $\beta_a=(z_a-\zeta)/(1-z_a\bar\zeta)$. From
$(1-z_a\bar\zeta)(1-z_b\zeta)-(z_a-\zeta)(z_b-\bar\zeta)=(1-z_az_b)(1-|\zeta|^2)$ we get
\[
 \frac{1-|\zeta|^2}{(1-z_a\bar\zeta)(1-z_b\zeta)}=\frac{1-\beta_a\bar\beta_b}{1-z_az_b}.
\]
Multiplying by $v_av_b(1-f_af_b)$ and summing gives
\begin{equation}\label{eq:phipsi}
 (1-|\zeta|^2)\bigl(|\psi(\zeta)|^2-|\varphi(\zeta)|^2\bigr)=u^*W_{\mathcal T}u,\qquad u_a=v_a\bar\beta_a,
\end{equation}
because $v^{\mathsf T}W_{\mathcal T}v=0$. So $|\varphi|\le|\psi|$ in the disc. For $|\zeta|=1$ one has
\[
 \frac1{(1-z_a\bar\zeta)(1-z_b\zeta)}=\frac1{1-z_az_b}\Bigl(\frac1{1-z_a\bar\zeta}+\frac1{1-z_b\zeta}-1\Bigr),
\]
and summation with $W_{\mathcal T}v=0$ gives $|\varphi|=|\psi|$ on the circle. Since the functions $1/(1-z_a\zeta)$ are independent, $\varphi\ne0$ and so $\psi\ne0$.
Hence $h=\varphi/\psi$ is a rational function with real coefficients, with $|h|\le1$ in the disc and $|h|=1$ on the circle. So it is a finite
Blaschke product. Both $\varphi$ and $\psi$ are polynomials of degree at most $|\mathcal T|-1$ divided by $\prod_a(1-z_a\zeta)$, so the degree of $h$ is at most
$|\mathcal T|-1$. Thus
\[
 h=\varepsilon z^s\frac{P}{P^*},\qquad \varepsilon=\pm1,\quad P^*(z)=z^{\deg P}P(1/z),
\]
where $P$ is a real polynomial with all roots in the open disc and $P(0)\ne0$. Put $d=s+\deg P$. Then $d+1\le k_++k_-$.

We have $h(z_b)=f_b$ for $b\in\mathcal T$. This is clear if $\psi(z_b)\ne0$. If $\psi(z_b)=0$, then $\varphi(z_b)=0$, and \eqref{eq:phipsi} at $\zeta=z_b$ gives $u^*W_{\mathcal T}u=0$
with $u_b=0$ and $u\ne0$. This contradicts the positive definiteness of $W_{\mathcal T\setminus\{b\}}$. We show that $h$ cannot exist.

\emph{The case $t\le s$.} The real polynomial $\varepsilon z^{s-t}P+cP^*$ is nonzero, because $|P|=|P^*|$ on the
unit circle and $c<1$. It vanishes at the $k_+$ positive nodes and has degree at most $d-t$. So $k_++t\le d\le
k_++k_--1$, that is, $t<k_-$. This contradicts \eqref{eq:window}, which gives $t\arctan4>\frac\pi2k_-$.

\emph{The case $s<t$.} Put $T=t-s\ge1$ and $r=h/z^t=\varepsilon P/(z^TP^*)$. Then $r=\gamma$ at the $k_-$ negative
nodes and $r=-c$ at the $k_+$ positive ones. On the unit circle $zh'/h=t+\Omega$, with $\Omega=zr'/r$ real.

\begin{lemma}\label{lem:contact}
$\Omega\ge0$ on the unit circle.
\end{lemma}

\begin{proof}
Let $v=\deg P$. The numerator of $zr'/r$ is the self-reciprocal real polynomial $zP'P^*-zP^{*\prime}P-TPP^*$ of
degree at most $2v$. So it equals $z^vR(z+1/z)$ with $\deg R\le v\le k_++k_--1$, and $\Omega=R(2\cos\theta)/|P|^2$
at $z=e^{i\theta}$. Rolle's theorem gives a critical point of $r$ between consecutive nodes of the same sign,
hence $k_++k_--2$ distinct roots of $R$ outside $[-2,2]$. So $R$ has at most one further root. Suppose that $R<0$
somewhere in $[-2,2]$. If $R(2)\ge0$ and $R(-2)\ge0$, then $R$ has two further roots in $[-2,2]$. If $R(2)<0$, then
$|r|$ decreases at $z=1$ on the real axis, so $|r|>1$ just to the left of $1$. Since $|r|<1$ at the largest node,
$r$ takes the value $r(1)=\pm1$ on the right of that node. Rolle's theorem gives one further critical point,
and a further root of $R$. The same holds at $-2$. If exactly one of $R(\pm2)$ is negative, this root comes with a sign
change of $R$ in $[-2,2]$. If both are negative, there are two such roots. In every case $R$ has two further roots,
which contradicts the degree bound.
\end{proof}

Let $p_1,\ldots,p_m$ be the positive nodes, so $m=k_+$ in the rest of this proof, and let $k=k_-$. We compare the
phase of $h$ with the phase of an explicit interpolant of degree $m$. Put $\lambda=c^2$ and $w_a=p_a^t$. Put
$\mathsf P_{ab}=(1-\lambda w_aw_b)/(1-p_ap_b)$ and $M=\mathsf P^{-1}$, and write $\langle a,b\rangle=a^{\mathsf T}Mb$. Let
$\xi_z(a)=1/(1-zp_a)$, $g_1=1$ and $g_2=-cw$. With $F_i(z)=g_i\xi_z$, $C_j=g_j\xi_{-1}$ and $J=\operatorname{diag}(1,-1)$ define
\[
 \Theta_{ij}(z)=\delta_{ij}-(1+z)J_j\langle F_i(z),C_j\rangle .
\]
So $\Theta(-1)=I$. Direct expansion gives the kernel identity
\begin{equation}\label{eq:kernelid}
 \sum_k\Theta_{ik}(z)J_k\Theta_{jk}(u)=\delta_{ij}J_i-(1-zu)\langle F_i(z),F_j(u)\rangle .
\end{equation}
The vector $F_1(p_a)-f_aF_2(p_a)$ is row $a$ of $\mathsf P$, where now $f_a=-cw_a$. This gives the node relation
$\Theta_{1j}(p_a)=f_a\Theta_{2j}(p_a)$. Let $D(z)=\prod_a(1-p_az)$, $b(z)=\prod_a(z-p_a)$ and
$\kappa=(-1)^m$. The functions $e_{ij}=D\Theta_{ij}$ are real polynomials of degree at most $m$. Their
determinant vanishes at the $p_a$ by the node relation and at the $1/p_a$ because $(e_{ij}(1/p_a))$ has rank one.
So it equals $\kappa bD$. Taking $u=1/z$ in \eqref{eq:kernelid} gives
\[
 e_{11}(z)=\kappa z^me_{22}(1/z),\qquad e_{12}(z)=\kappa z^me_{21}(1/z).
\]
On the unit circle \eqref{eq:kernelid} with $i=j=2$ and $u=\bar z$ gives $|\Theta_{22}|^2=1+|\Theta_{21}|^2$. Let
$\sigma=h(-1)=\pm1$ and $e=\sigma e_{21}+e_{22}$. Then $W=e/D=\Theta_{22}(1+\sigma\Theta_{21}/\Theta_{22})$ does not
vanish on the circle, and $W(-1)=1$.

\emph{The size of the argument of $\Theta_{22}$.} Let $T_l$ be $\Theta_{22}$ with $\lambda$ replaced by $l$, so
$T_0=1$. Let $\cK^{(p)}$ be the matrix with entries $1/(1-p_ap_b)$, let $Z=\operatorname{diag}(w_a)$, and let
$M_l=(\cK^{(p)}-lZ\cK^{(p)}Z)^{-1}$. For $0\le l<l'<1$ and real vectors $a,b$ the resolvent identity gives
\[
 l'\,a^{\mathsf T}ZM_{l'}Zb-l\,a^{\mathsf T}ZM_lZb=(l'-l)(M_{l'}Za)^{\mathsf T}\cK^{(p)}(M_lZb) .
\]
Since $\cK^{(p)}-Z\cK^{(p)}Z\ge0$, the right side has modulus at most $(l'-l)\|a\|\|b\|/((1-l)(1-l'))$, where
$\|a\|^2=a^{\mathsf T}(\cK^{(p)})^{-1}a$. For $c=0$ the identity \eqref{eq:kernelid} at $z=\zeta$ and $u=\bar\zeta$ gives
$(1-|\zeta|^2)\|\xi_\zeta\|^2=1-|b(\zeta)/D(\zeta)|^2$. Since $1-\prod t_a\le\sum(1-t_a)$ for $t_a\in[0,1]$, this yields
\[
 \|\xi_\zeta\|^2\le\sum_a\frac{1-p_a^2}{|1-p_a\zeta|^2}\qquad(|\zeta|\le1).
\]
In particular $\|\xi_{-1}\|^2\le m$. For $z$ on the arc from $z_0$ to $-1$ one has $\operatorname{Re}z\le\frac{15}{17}$,
and a quadratic inequality in $p$ gives $|1+z|^2(1-p^2)/|1-pz|^2\le(\frac{17}4)^2$. So
$|1+z|\,\|\xi_z\|\,\|\xi_{-1}\|\le\frac{17}4m$, and
$|T_{l'}(z)-T_l(z)|\le\frac{l'-l}{(1-l)(1-l')}\cdot\frac{17}4m$ on the arc. Since $|T_l|\ge1$, a telescoping product
over a fine grid of $[0,\lambda]$ shows that the continuous argument of $\Theta_{22}$ on the arc, normalized to vanish
at $-1$, has modulus at most $\mathcal X=\frac{17}4m\lambda/(1-\lambda)$. The factor $1+\sigma\Theta_{21}/\Theta_{22}$ has
positive real part. So the continuous argument $\ell(\theta)$ of $W(e^{i\theta})$, with $\ell(\pi)=0$, satisfies
$|\ell(\theta_0)|\le\frac\pi2+\mathcal X$, where $z_0=e^{i\theta_0}$ and $\pi-\theta_0=2\arctan4$.

\emph{The phases.} By Lemma~\ref{lem:contact}, $h(e^{i\theta})=\sigma e^{-i\phi(\theta)}$ with
$\phi(\theta)=\int_\theta^\pi(t+\Omega)\ge t(\pi-\theta)$. Let
$\psi(\theta)=m(\pi-\theta)+2\sum_a\arg(1-p_ae^{i\theta})$. Then $(-z)^m\bar e/e=\exp(-i(\psi+2\ell))$. At $\theta_0$,
$1-pz_0=(1+p)(1+4i\rho_p)/(1+4i)$. So $\psi(\theta_0)=2\sum_p\arctan(4\rho_p)$.

\emph{The root count.} Let $\hat e(z)=z^me(1/z)$ and
\[
 \mathcal N=\varepsilon z^sPe-\sigma\kappa P^*\hat e,
\]
a real polynomial of degree at most $d+m$. Put $a=\phi-\psi-2\ell$, so $a(\pi)=0$. At every $\theta$ with
$a(\theta)\in2\pi\Z$ the polynomial $\mathcal N$ vanishes at $e^{\pm i\theta}$. It also vanishes at every $p_a$ by the
node relation and the reversal formulas, at every $1/p_a$ because $z^{d+m}\mathcal N(1/z)=\pm\mathcal N(z)$, and at $-1$. Suppose
that $a(\theta_0)>\pi k$ and let $L=\lfloor k/2\rfloor$. If $\mathcal N\ne0$, the intermediate value theorem gives $L$
conjugate pairs of roots on the open arc. So $2L+2m+1\le d+m\le2m+k-1$, which contradicts $2L+2>k$. If
$\mathcal N=0$, then $e^{ia}=1$ on the arc, which contradicts the intermediate value theorem at the value $\pi$.
Hence $a(\theta_0)\le\pi k$, and
\[
 2t\arctan4\le\phi(\theta_0)\le\pi k+2\sum_p\arctan(4\rho_p)+\pi+2\mathcal X .
\]
This contradicts \eqref{eq:window}, and Theorem~\ref{thm:window} is proved.\qed

\subsection{The alternative on near-extremal configurations}\label{ssec:alternative}
We now show that \eqref{eq:window} holds with $t=\mathsf e$ and $c=\frac1{100}$ on every configuration for which the
energy bound of Section~\ref{sec:energy} does not already suffice. The idea is that such a configuration has
Chebyshev moments close to the trial sequence $u$, and the left side of \eqref{eq:window} is then computed by a
fixed test function.

Let $\phi_H(w)=\frac\pi2-\arctan\bigl(4(1-w)/(1+w)\bigr)$ for $0\le w<1$. Then $\arctan(4\rho_p)=\frac\pi2-\phi_H(p)$.
With $\mathsf e=48N-1$, $n=58N$ and $k_-=n-k_+$, the inequality \eqref{eq:window} follows for large $N$ from
\begin{equation}\label{eq:phaseneed}
 \frac1n\sum_{p>0}\phi_H(p)\ge\arctan\tfrac14+\tfrac5{29}\arctan4+\tfrac{17}4\,\frac{c^2}{1-c^2}+\xi,\qquad \xi=\tfrac1{1000},
\end{equation}
whose right side is $0.474992\ldots$. Let $f=\frac{F_0}2+\sum_{k\ge1}F_kT_k$ be a polynomial of degree $32$ with
dyadic coefficients and denominator $2^{31}$, where $F_0=992223195/2^{30}$. 
It satisfies $f\le\phi_H$ on $(0,1)$, $f\le0$ on $[-1,0]$, and
\begin{equation}\label{eq:testS}
 \tfrac14F_0+\tfrac14|F_1|y\le\phi_H(y)\qquad(0<y<1).
\end{equation}
Also $\sum_kkF_k^2\le7877/10000$. On this range of $\alpha$ the case with doubled interaction uses the moving trial
sequence $u(\tau)=u+\tau\,\mathrm du$ of \S\ref{ssec:trials}, with the increment of $40$ Chebyshev coefficients and
$\tau=\alpha/\alpha_1$. We have $\sum_kF_ku_k\ge879/10000$ and $\sum_kF_k\mathrm du_k\ge-0.010526$.

Fix a configuration and a set $S$. Let $\mu_F$, $\mu_M$ and $\mu_S$ be the measures of \S\ref{ssec:prescribed}, with the points of $S$ as the
prescribed points of the second kind, and let $\beta$ be the mass of $\mu_S$. By \eqref{eq:squares}, the bound of Lemma~\ref{lem:energy-ineq} with
$\alpha_2=\beta$ improves by the number $\Delta$ of \eqref{eq:squares}. Divide $[0,\alpha_1]$ into $64$ equal cells. On the cell $j$ the
left side of the inequality of Lemma~\ref{lem:energy-ineq}, minus the right side at the slope $c$, is at most a rational number $\Delta_j$.
These numbers lie between $10^{-6}$ and $0.008365$. So either $\Delta\le\Delta_j$, or the energy bound gives the conclusion of
Proposition~\ref{prop:signed-goal} for $\Xi_S$ directly.

Suppose that $\Delta\le\Delta_j$. Since $\phi_H\ge0$ on $(0,1)$ and $f\le\phi_H$, while $f\le0$ on the negative
half, we have $\frac12f\le\phi_H$ on positive prescribed points and $\frac12f\le0$ on negative ones. With
\eqref{eq:testS},
\begin{align*}
 \frac1n\sum_{p>0}\phi_H(p)&\ge\int f\,d(\mu_F+\tfrac12\mu_M)+\tfrac14F_0\beta+\tfrac14|F_1|A_1(\mu_S)\\
 &\ge\tfrac12F_0-\tfrac14F_0\alpha+\sum_kF_ku_k(\tau)-\sqrt{\tfrac12\Delta_j\sum_kkF_k^2},
\end{align*}
by the Cauchy--Schwarz inequality. The right side is affine in $\alpha$ on each cell. At the end points of every cell it exceeds
$0.4752$, which is larger than the right side of \eqref{eq:phaseneed}. The smallest margin is $0.0050$. So \eqref{eq:phaseneed} holds, and
Theorem~\ref{thm:window} gives $F(-\frac1{100})\ge0$. All the numerical facts of this subsection are part of the formal proof, which also carries the regularization of the moments that we have suppressed here. We have proved the following.

\begin{proposition}\label{prop:alternative}
For $N$ large, every configuration with $m\le\alpha_1n$ prescribed columns in the case with doubled interaction, and
every $S$, satisfies at least one of two statements. Either the energy bound for $\Xi_S$ gives the conclusion
of Proposition~\ref{prop:signed-goal}, or $F(-\frac1{100})\ge0$ and therefore
$|\det\Xi_S|\le499^n|\det\Xi^0_S|$.
\end{proposition}

\subsection{Gram determinants}\label{ssec:gram}
It remains to bound $\det\Xi^0_S$. Let $\mathsf D=n-1-g=53N-1$. Let $\xi_1,\ldots,\xi_{k_+}$ be the positive nodes, with
$\xi=z$ outside $S$ and $\xi=1/z$ on $S$, so that all these rows are $(\xi^{g-r})_r$. Let $x_1,\ldots,x_{k_-}$ be the
negative nodes, and put $u_b=J(x_b)$, where $J(z)=\frac12(z+1/z)$. Let $s=\mathsf D+1-k_-$ and $\omega(z)=\prod_p(z-\xi_p)$. If
$k_->\mathsf D+1$, the determinant vanishes, because the negative rows lie in a space of dimension $\mathsf D+1$. So we may assume
$s\ge0$. The Laurent polynomials with exponents in $[-\mathsf D,g]$ that vanish at all $\xi_p$ are spanned by the functions
$z^{i-\mathsf D}\omega(z)$ with $0\le i<k_-$. So elimination gives
\begin{equation}\label{eq:split}
 |\det\Xi^0_S|=\prod_p|\xi_p|^{-\mathsf D}\,|V(\xi)|\,\bigl|\det\bigl(\mathcal Fz^{i-\mathsf D}\omega\,(x_b)\bigr)_{i,b}\bigr|,
\end{equation}
where $\mathcal Ff(x)=\frac12(f(x)+f(1/x))$ is the fold. The fold sends $z^d$ to $T_{|d|}(J(z))$. So the fold of a Laurent
polynomial of degree at most $\mathsf D$ is a polynomial of degree at most $\mathsf D$ in $u=J(z)$.

We measure these functions in a finite Laurent-orthogonal measure. Let $\nu=\nu_N$ be the positive measure on $16$
concentric circles of radii $r_j=e^{-\sigma_j}$, with total masses $e^{-N\tau_j}$, which is uniform on the roots of unity
of order $2\mathsf D+1$ on each circle (scaled by $r_j$). We assume that $\sigma_0=\tau_0=0$ and $10\sigma_j\le\tau_j$ for all
$j$. The Laurent monomials with exponents in $[-\mathsf D,\mathsf D]$ are orthogonal in $L^2(\nu)$, with $\|z^c\|^2=A_c=\sum_je^{-N\tau_j}r_j^{2c}$. Let $G_0=A_0$ and $G_a=(A_a+A_{-a})/4$ for $a\ge1$, which is the squared
norm of $T_a(J(z))$.

\begin{proposition}\label{prop:gram}
Let $\ell_a$ be the leading coefficient of $T_a$, and suppose that the numbers $\theta_d\ge1$ for $-\mathsf D\le d\le g$ satisfy
$G_a\sum_{d}(A_d\theta_d)^{-1}\le1$ for each $a\le \mathsf D$, the sum running over the $d$ with $|d|=a$. Then
\begin{multline*}
 (\det\Xi^0_S)^2\prod_{a=0}^DG_a\cdot k_-!\,s!\\
 \le\Bigl(\prod_p\xi_p^{-\mathsf D}V(\xi)\Bigr)^2\Bigl(\prod_{a=0}^{\mathsf D}\ell_a\,V(u)\Bigr)^2\\
 \cdot\prod_d\theta_d\cdot Z_+Z_-,
\end{multline*}
where
\begin{gather*}
 Z_+=\int\prod_i|z_i|^{-2\mathsf D}|\omega(z_i)|^2|V(z)|^2\,d\nu^{k_-},\\
 Z_-=\int\prod_l|\pi(Jz_l)|^2|V(Jz)|^2\,d\nu^{s},
\end{gather*}
with $\pi(u)=\prod_b(u-u_b)$.
\end{proposition}

\begin{proof}
Let $\mathcal P$ be the space of polynomials of degree at most $\mathsf D$ in $u$, with the inner product for which the $T_a$ are
orthogonal of squared norm $G_a$. Consider the two families $f_i=\mathcal Fz^{i-\mathsf D}\omega$ for $0\le i<k_-$ and $\pi u^l$ for $0\le l<s$.
They have $\mathsf D+1$ members together. The second family spans the kernel of the evaluation at the $u_b$. So the
determinant of the coefficients of the joint family in the basis $(T_a)$ equals $\det(f_i(x_b))$ divided by
$\prod_a\ell_a\cdot V(u)$, up to sign. Multiplying its square by $\prod_aG_a$ gives the Gram determinant of the joint family.
Fischer's inequality bounds it by the product of the Gram determinants of the two families.

Let $f=\sum_df_dz^d$ be a combination of the functions $z^{i-\mathsf D}\omega$. The coefficient of $T_a$ in $\mathcal Ff$ is
$\sum_{|d|=a}f_d$, and the Cauchy--Schwarz inequality with the hypothesis gives $\|\mathcal Ff\|^2\le\sum_d\theta_dA_df_d^2$. Raising
the weights of a diagonal form by factors $\theta_d\ge1$ raises a Gram determinant by at most $\prod_d\theta_d$, by the
Cauchy--Binet formula. So the first Gram determinant is at most $\prod_d\theta_d$ times the Gram determinant of the
functions $z^{i-\mathsf D}\omega$ in $L^2(\nu)$. Andr\'eief's identity \cite{Forrester2018Andreief} in its finite form writes the two Gram
determinants as $Z_+/k_-!$ and $Z_-/s!$. The inequality follows with \eqref{eq:split}.
\end{proof}

\emph{The constants $\theta_d$ and $G_a$.} Put $\theta_d=1$ for $d\le0$ and $\theta_d=A_{-d}/A_d$ for $1\le d\le g$. For $a>g$ the
sum in the hypothesis has the single term $d=-a$, and $G_a/A_{-a}\le\frac12$ because $A_a\le A_{-a}$. For $1\le a\le g$
the left side of the hypothesis is $\frac12(A_a+A_{-a})/A_{-a}\le1$. Since $g=5N$ and $10\sigma_j\le\tau_j$, each term of $A_{-a}$ is at
most $1$, so $A_{-a}\le16$ for $a\le g$. Also $A_a\ge1$ by the term $j=0$. Hence $\prod_d\theta_d\le16^{5N}$. Let
$\tau^*(q)=\max_j(2q\sigma_j-\tau_j)$. Then $G_a\ge\frac14e^{N\tau^*(a/N)}$ for $a\ge1$. Since $\tau^*$ is nondecreasing,
\begin{equation}\label{eq:legendre}
 \sum_{a=0}^{\mathsf D}\log G_a\ge N^2\int_0^{53}\tau^*(q)\,dq-2\mathsf D\log2-N\tau^*(53).
\end{equation}
The factor $\prod_a\ell_a^2$ equals $2^{\mathsf D(\mathsf D-1)}$.

\emph{A tangent bound for Vandermonde determinants.} The integrals $Z_\pm$ are compared with reference measures by the
following lemma, which is the negative definiteness of the logarithmic energy on neutral measures in a form exact
for finitely many points. Let $d(z,w)^2=|z-w|^2/((1+|z|^2)(1+|w|^2))$ and $\mathsf h(z)=\frac12\log(1+|z|^2)$. For a finite positive measure $q$ on
$\C$ write $U^q(z)=\int\log|z-w|\,dq(w)$ and $I(q)=\int U^q\,dq$. Let $H_L=\sum_{j\le L}1/j$.

\begin{lemma}\label{lem:tangent}
Let $q$ be a finite positive measure on $\C$ with $\int\mathsf h\,dq<\infty$ and $q\{w:d(z,w)^2\le t^4\}\le Ct$ for all $z$ and all
$0<t\le1$. For every $\eta>0$ there is $L$ such that for every $m\ge0$, all $\nu_0\ge0$ with $\nu_0q(\C)=m$ and all
$z_1,\ldots,z_m\in\C$,
\[
 |V(z)|^2\le\exp\Bigl(2\nu_0\sum_iU^q(z_i)-\nu_0^2I(q)+2\nu_0m\eta+\tfrac m2H_L\Bigr).
\]
\end{lemma}

\begin{proof}
We may assume that the $z_i$ are distinct. Let $\mathsf n(z)$ be the stereographic image of $z$ on the unit sphere. Then
$\mathsf C(z,w)=1-d(z,w)^2=\frac12(1+\mathsf n(z)\cdot\mathsf n(w))$ is a sum of four products $b_if_i(z)f_i(w)$ with $b_i\ge0$, and so
is every power of $\mathsf C$, with more terms. Let $K_L=-\frac12\sum_{j\le L}\mathsf C^j/j$. Then $K_L=-\sum_\phi b_\phi F_\phi\otimes F_\phi$ is a finite sum with
$b_\phi\ge0$. For real $S_\phi=\sum_iF_\phi(z_i)$ and $M_\phi=\int F_\phi\,dq$ the inequality $(S_\phi-\nu_0M_\phi)^2\ge0$ gives, after summation,
\begin{equation}\label{eq:gramtangent}
 \sum_{i,j}K_L(z_i,z_j)\le2\nu_0\sum_i\int K_L(z_i,w)\,dq(w)-\nu_0^2\iint K_L\,dq\,dq .
\end{equation}
For $z\ne w$ we have $\log d(z,w)=\frac12\log(1-\mathsf C)=-\frac12\sum_{j\ge1}\mathsf C^j/j\le K_L(z,w)$, and $K_L(z,z)=-\frac12H_L$. So
$\sum_{i\ne j}\log d(z_i,z_j)$ is at most the left side of \eqref{eq:gramtangent} plus $\frac m2H_L$. Also $\iint K_L\ge\iint\log d$. For the
remaining term we show that $\int K_L(z,\cdot)\,dq\le\int\log d(z,\cdot)\,dq+\eta$ for large $L$, uniformly in $z$. Fix $h\in(0,1)$.
On $\{d^2\ge h\}$ the difference $K_L-\log d=\frac12\sum_{j>L}\mathsf C^j/j$ is at most $(1-h)^L/(2Lh)$. On $\{d^2<h\}$ it is at most $-\log d$, and
the ball hypothesis on the shells $2^{-4(i+1)}<d^2\le2^{-4i}$ bounds the integral of $-\log d$ over this set by a quantity that tends to
$0$ with $h$, uniformly in $z$. This proves the claim.

Finally $\log|z-w|=\log d(z,w)+\mathsf h(z)+\mathsf h(w)$. So $\sum_{i\ne j}\log|z_i-z_j|=\sum_{i\ne j}\log d(z_i,z_j)+2(m-1)\sum_i\mathsf h(z_i)$. Expanding
$U^q$ and $I(q)$ in the same way, and using $\nu_0q(\C)=m$, the terms in $\int\mathsf h\,dq$ cancel. What remains is
$2m\sum_i\mathsf h(z_i)$ on the right side against $2(m-1)\sum_i\mathsf h(z_i)$ on the left side. Since $\mathsf h\ge0$, this proves the lemma.
\end{proof}

Apply the lemma to the $a+c$ points consisting of $a$ fixed base points $b_i$ and $c$ points $Z_l$. If $Z_l$ runs over a finite set with
weights $\mathrm{wt}\ge0$, summing over tuples gives, with $\nu_0q(\C)=a+c$ and the same $\eta$ and $L$,
\begin{multline}\label{eq:tangent}
 |V(b)|^2\sum_{Z}\prod_l\mathrm{wt}(Z_l)\,\prod_{l,i}|Z_l-b_i|^2\,|V(Z)|^2\\
 \le\exp\Bigl(2\nu_0\sum_iU^q(b_i)-\nu_0^2I(q)+2\nu_0(a+c)\eta+\tfrac{a+c}2H_L\Bigr)\Bigl(\sum_z\mathrm{wt}(z)e^{2\nu_0U^q(z)}\Bigr)^c .
\end{multline}
The exponential on the right does not depend on the $Z_l$, and the last sum separates them. This is the reason for completing the
gas by the base points: the determinant of the base points is then bounded together with the gas.

\emph{The two gases.} Let $q_+$ be a measure of mass $1$ in the plane of $z$, and $q_-$ a measure of mass $\frac{53}{58}$ in the plane of
$u$, both satisfying the hypotheses of Lemma~\ref{lem:tangent}. Let $\kappa_+$ and $c_-$ be constants. We assume the angular bounds
\begin{equation}\label{eq:angular}
 U^{q_+}(r_je^{i\theta})+\tfrac{53}{58}\sigma_j-\tfrac{\tau_j}{116}\le\kappa_+,\qquad
 U^{q_-}\bigl(J(r_je^{i\theta})\bigr)-\tfrac{\tau_j}{116}\le c_-,
\end{equation}
for all $j$ and $\theta$. Apply \eqref{eq:tangent} to $Z_+$ with $q=q_+$, $\nu_0=n$, base points $\xi_p$ and the $k_-$ gas points $z_i$. Since
$r_j^{-2\mathsf D}\le e^{2\cdot53N\sigma_j}$, the sum in \eqref{eq:tangent} is at most $16e^{2n\kappa_+}$. Apply \eqref{eq:tangent} to $Z_-$ with $q=q_-$,
$\nu_0=n$, base points $u_b$ and the $s$ gas points $Jz_l$. The sum is at most $16e^{2nc_-}$. Insert these in
Proposition~\ref{prop:gram}, use \eqref{eq:legendre}, and take logarithms. Each positive node contributes
$\mathsf g(\xi)=-\mathsf D\log|\xi|+n(U^{q_+}(\xi)-\kappa_+)$ and each negative node $\mathsf g(x)=n(U^{q_-}(Jx)-c_-)$. We have proved the following statement.

\begin{proposition}\label{prop:gas}
For every $\eta>0$ there is $C$ such that for all $N$, all nodes and all $S$ with $\det\Xi^0_S\ne0$,
\[
 2\log|\det\Xi^0_S|\le2\sum_a\mathsf g(a)+\mathsf K_N+4n^2\eta+Cn,
\]
where $\mathsf g(a)$ is the contribution of node $a$ and
\[
 \mathsf K_N=2n^2\kappa_++2\cdot53Nn\,c_-+\mathsf D(\mathsf D+1)\log2-N^2\int_0^{53}\tau^*-n^2\bigl(I(q_+)+I(q_-)\bigr).
\]
\end{proposition}

The constant is correct because the $k_-$ points of $Z_+$ give $2nk_-\kappa_+$ and the $s$ points of $Z_-$ give $2nsc_-$. In $\mathsf K_N$ these are
written as $2n^2\kappa_+$ and $2\cdot53Nnc_-$, and the differences are absorbed by the terms $-n\kappa_+$ and $-nc_-$ of $\mathsf g$. The term $\mathsf D(\mathsf D+1)\log2$
is the sum of $\mathsf D(\mathsf D-1)\log2$ from $\prod_a\ell_a^2$ and $2\mathsf D\log2$ from \eqref{eq:legendre}.

\subsection{The final inequality}\label{ssec:budget}
We now bound the other factors of the integrand. Let $n'=n-m$, with free points $x_i$ and $s_j$ and prescribed points $y_j$, and let $\mathcal S$ be the product of the
factors of the integrand other than the row determinant and the factors $\mathcal B_\zeta$ of Lemma~\ref{lem:binomvand}. Put $\mathsf t(x)=2x/(1+x^2)$ and $t_i=\mathsf t(x_i)$. In the normalization of
\S\ref{ssec:rowdet} a free point carries the weight $|t|^{75N-1}(1-t)^{3N}$. The prescribed columns have lost the common factor $|t|^{22N}(1-t)^{3N}$ of the rows, so a prescribed point
carries $|\mathsf t(y)|^{\mathsf D}$. The Cauchy determinant of Lemma~\ref{lem:andreief} is $\pm V(t)V(s)/\prod_{i,j}(1-t_is_j)$, and the generalized Vandermonde determinant gives
$\prod_js_j^{9N}(1-s_j)^{3N}|V(s)|$. Hence
\begin{multline*}
 \mathcal S=\prod_j|\mathsf t(y_j)|^{\mathsf D}\prod_i|t_i|^{75N-1}(1-t_i)^{3N}\\
 \cdot\prod_js_j^{9N}(1-s_j)^{3N}\,\frac{|V(t)|\,V(s)^2}{\prod_{i,j}(1-t_is_j)}.
\end{multline*}
Now $1-t=(1-x)^2/(1+x^2)$, $1-ts=(1-2xs+x^2)/(1+x^2)$ and
\[
 t_i-t_{i'}=\frac{2(x_i-x_{i'})(1-x_ix_{i'})}{(1+x_i^2)(1+x_{i'}^2)}.
\]
Each $x_i$ occurs in $n'-1$ differences and in $n'$ kernels, so the exponent of $1+x_i^2$ is $-(75N-1)-3N-(n'-1)+n'=-(78N-2)$. Taking logarithms gives
\begin{multline}\label{eq:scalar}
 \log\mathcal S=\Bigl(n'(75N-1)+\binom{n'}2\Bigr)\log2+\sum_j\mathsf w(y_j)+\sum_i\mathsf W_N(x_i)\\
 +\sum_j\bigl(9N\log s_j+3N\log(1-s_j)\bigr)+n^2\mathsf Q(x,s),
\end{multline}
where $\mathsf w(y)=\mathsf D(\log2+\log|y|-\log(1+y^2))$, $\mathsf W_N(x)=(75N-1)\log|x|+6N\log(1-x)-(78N-2)\log(1+x^2)$ and
\begin{multline*}
 n^2\mathsf Q(x,s)=\sum_{i<i'}\log\bigl(|x_i-x_{i'}|(1-x_ix_{i'})\bigr)\\
 +2\sum_{j<j'}\log|s_j-s_{j'}|-\sum_{i,j}\log(1-2x_is_j+x_i^2).
\end{multline*}
This is the energy of \S\ref{ssec:energy} with the interaction among the $x_i$ at weight $1$. The second copy of that interaction is
contained in the determinant.

For sequences $a=(a_k)$ and $v=(v_k)$ put $\|a\|^2=\sum_ka_k^2/k$. Let $b$ be a real number.

\begin{lemma}\label{lem:freetangent}
Let $a,v$ be sequences with $\sum|a_k|/k$, $\sum|v_k|/k$, $\|a\|$ and $\|v\|$ finite. For every $\varepsilon>0$ there is $\nu_1$ such that
for all $n\ge\nu_1$, all $n'\le n$ and all configurations of distinct points,
\begin{multline*}
 \mathsf Q(x,s)\le\tfrac32b^2\log2+\|a\|^2+\tfrac12\|v\|^2-\tfrac32\bigl(\tfrac{n'}n-b\bigr)^2\log2\\
 +\frac1n\sum_i\mathsf R_*(x_i)+\frac1n\sum_j\mathsf C_*(s_j)\\
 -\frac1{2n^2}\sum_i\log(1-x_i)-\frac1{n^2}\sum_j\log(1-s_j)+\varepsilon,
\end{multline*}
where $\mathsf R_*(x)=-b\log2-2\mathsf T(a,x)-\mathsf S(v,x)$ and $\mathsf C_*(s)=-2b\log2+2\mathsf T(v,s)$.
\end{lemma}

\begin{proof}
Use the expansions $-\log|x-x'|=\log2+2\sum_kT_k(x)T_k(x')/k$, $-\log(1-xx')=\sum_kx^kx'^k/k$ and
$-\log(1-2xs+x^2)=2\sum_kx^kT_k(s)/k$. Let $A_k=\frac1n\sum_iT_k(x_i)$, $B_k=\frac1n\sum_ix_i^k$ and $C_k=\frac1n\sum_jT_k(s_j)$, and let $f=n'/n$.
If the diagonal terms are ignored, these expansions give
\[
 \mathsf Q=-\tfrac32f^2\log2-\|A\|^2-\tfrac12\|B-2C\|^2 .
\]
The diagonal terms diverge. OpenAI's regularization \cite[\S5]{OpenAI2026} replaces the kernels by Abel-damped series with a parameter $e$.
It gives the identity above for the damped moments, up to an error at most
$-\frac1{2n^2}\sum_i\log(1-x_i)-\frac1{n^2}\sum_j\log(1-s_j)+O_e(1/n)+14e$. For the trial sequences, $-\|A\|^2\le\|a\|^2-2\sum_ka_kA_k/k$ and
$-\frac12\|B-2C\|^2\le\frac12\|v\|^2-\sum_kv_k(B_k-2C_k)/k$. Removing the damping from the linear terms costs $O(e)$ times the $\ell^1$ norms of the trial
sequences, so we first replace $a$ and $v$ by long prefixes, with an error at most $\varepsilon/4$ uniformly in the points. Finally
$-\frac32f^2\log2=-\frac32(f-b)^2\log2+\frac32b^2\log2-3bf\log2$, and the last term is split between $\mathsf R_*$ and $\mathsf C_*$. Choosing the prefix, then $e$, then $\nu_1$
proves the lemma. The order of these choices makes the bound uniform in the configuration.
\end{proof}

We apply the lemma with $b=1$. Let $X$ be the inverse of $x\mapsto x+1/x$ with values in $[-1,1]\setminus\{0\}$. Let $R$ and $C$ be the row and column references of Appendix~\ref{app:trials},
and let $L(\mu)=\int\log|X|\,d\mu$. We take
\[
 a_k=\int T_k(X)\,dR,\qquad v_k=\int X^k\,dR-2\int T_k\,dC,
\]
and put $\mathsf I_0=\frac32\log2+\|a\|^2+\frac12\|v\|^2$. The expansions in the proof of Lemma~\ref{lem:freetangent} give
\begin{gather*}
 \mathsf R_*(x)=\int\bigl(\log|x-X|+\log(1-xX)\bigr)\,dR-\int\log(1-2xs+x^2)\,dC(s),\\
 \mathsf C_*(s)=2U^C(s)-\int\log(1-2Xs+X^2)\,dR.
\end{gather*}
For nonzero $x$ and $w$ and for $0<s<1$ we have
\begin{equation}\label{eq:jouk}
 x+\frac1x-w-\frac1w=\frac{(x-w)(xw-1)}{xw},\qquad 1-2xs+x^2=|x|\,\bigl|x+\tfrac1x-2s\bigr|.
\end{equation}
Since $R$ and $C$ have mass $1$, this gives
\[
 \mathsf R_*(x)=L(R)+U^{R-2_\#C}(x+1/x),\qquad \mathsf C_*(s)=-L(R)+2U^C(s)-U^R(2s).
\]
So $9N\log s+3N\log(1-s)+n\mathsf C_*(s)=n\mathsf C(s)$ for the column field $\mathsf C$ of the appendix.

\emph{The moved sequences.} Lemma~\ref{lem:freetangent} holds for any sequences, so we may move $a$ and $v$ with the
fraction $\alpha$. Let $\mathrm da^0,\mathrm da^1,\mathrm dv^0,\mathrm dv^1$ be fixed sequences supported on $k\le160$, with
dyadic entries of denominator $2^{40}$. Put $\alpha_e=\frac1{25}$, $\tau=\min(1,\alpha/\alpha_e)$ and
\[
 a(\tau)=a+\mathrm da^0+\tau\,\mathrm da^1,\qquad v(\tau)=v+\mathrm dv^0+\tau\,\mathrm dv^1 .
\]
The increment $(\mathrm da^0,\mathrm dv^0)$ moves the sequences already at $\alpha=0$. The fields change by finite sums:
$\mathsf R_*$ gains $h^0+\tau h^1$ with $h^i=-2\mathsf T(\mathrm da^i,\cdot)-\mathsf S(\mathrm dv^i,\cdot)$, and $\mathsf C_*$ gains
$k^0+\tau k^1$ with $k^i=2\mathsf T(\mathrm dv^i,\cdot)$. The polarization identity gives
\[
 \mathsf I_0(\tau)=\mathsf I_0+(L_0+Q_{00})+\tau(L_1+2Q_{01})+\tau^2Q_{11},
\]
with $L_i=2\langle a,\mathrm da^i\rangle+\langle v,\mathrm dv^i\rangle$ and
$Q_{ij}=\langle\mathrm da^i,\mathrm da^j\rangle+\frac12\langle\mathrm dv^i,\mathrm dv^j\rangle$. Only $L_0$ and $L_1$ involve the
references. We also include the column field $\gamma g$ of \S\ref{ssec:prescribed}, with $\gamma=\tau\alpha_e$ for $\alpha\le\alpha_e$
and $\gamma=\alpha$ for $\alpha_e\le\alpha\le\alpha_1$, so that $\gamma\ge\alpha$. The moved row field at $\tau$ is the convex
combination $(1-\tau)(F_\pm+h^0)+\tau(F_\pm+h^0+h^1)$, and the moved column field is affine in $\tau$ and $\gamma$. So five
numbers bound all of them, the bounds at the end points:
\begin{gather*}
 F_\pm+h^0\le c^{\mathsf R}_0,\qquad F_\pm+h^0+h^1\le c^{\mathsf R}_1,\\
 \mathsf C+k^0\le c^{\mathsf C}_0,\qquad \mathsf C+k^0+k^1+\alpha_eg\le c^{\mathsf C}_1,\qquad \mathsf C+k^0+k^1+\alpha_1g\le c^{\mathsf C}_2 .
\end{gather*}
For $\alpha\le\alpha_e$ the row and column caps are the chords $(1-\tau)c^{\mathsf R}_0+\tau c^{\mathsf R}_1$ and
$(1-\tau)c^{\mathsf C}_0+\tau c^{\mathsf C}_1$. For $\alpha_e\le\alpha\le\alpha_1$ they are $c^{\mathsf R}_1$ and the chord of
$c^{\mathsf C}_1$ and $c^{\mathsf C}_2$ at $\alpha$. No supremum of an increment is ever taken separately from the supremum of the base field.

We turn to the gas potentials. Let $\nu_\pm$, $A_\pm$ and $\tilde q_-$ be as in Appendix~\ref{app:trials}, so that $q_+=X_\#G_++\nu_+$ and $q_-=J_\#\tilde q_-$ with $\tilde q_-=X_\#G_-+\nu_-$.
Let $0<x<1$ and $Y=x+1/x$. By \eqref{eq:jouk}, $\log|x-X(Y')|=\log|Y-Y'|+\log x+\log|X(Y')|-\log(1-xX(Y'))$. On the support of $G_+$ we have $X>0$, so
$-\log(1-xX)=\sum_kx^kX^k/k$ has positive terms. Those with $k\le900$ integrate to $\mathsf P(x)$ and the others to at most $\mathsf E$. Hence
\begin{equation}\label{eq:qplus}
 U^{q_+}(x)\le L(G_+)+U^{G_+}(Y)+m_+\log x+\mathsf P(x)+\mathsf E+A_+(x).
\end{equation}
Let now $-1<x<0$. We have $\log|Jx-Y'/2|=\log|Y-Y'|-\log2$ on the support of $G_-$. By \eqref{eq:jouk},
$\log|Jx-Jw|=\log|x-w|+\log|1-xw|-\log2-\log|x|-\log|w|$ on the support of $\nu_-$. Put $c_M=\int\log|u|\,d\tilde q_-(u)$. Integrating, and using the mass $\frac{53}{58}$ of $\tilde q_-$, we find
\begin{multline}\label{eq:qminus}
 U^{q_-}(Jx)=L(G_-)+U^{G_-}(Y)+m_-\log|x|+A_-(x)\\
 -\tfrac{53}{58}\bigl(\log2+\log|x|\bigr)-c_M.
\end{multline}
Define $\kappa_-$ by $c_-=\kappa_--\frac{53}{58}\log2-c_M$. Since $n\cdot\frac{53}{58}=53N$, this says
\begin{equation}\label{eq:gneg}
 \mathsf g(x)=n\bigl(L(G_-)+U^{G_-}(Y)+m_-\log|x|+A_-(x)-\kappa_-\bigr)-53N\log|x|.
\end{equation}

Let $F_\pm$ be the free row fields of the appendix. Since $n\cdot\frac{11}{29}=22N$, $n\cdot\frac{39}{29}=78N$ and $n\cdot\frac3{29}=6N$, the definition reads
\begin{multline*}
 nF_\sigma=n\mathsf R_*+n\bigl(L(G_\sigma)+U^{G_\sigma}(Y)+m_\sigma\log|x|-\kappa_\sigma\\
 +[\sigma=+](\mathsf P+\mathsf E)+A_\sigma\bigr)+22N\log|x|-78N\log(1+x^2)+6N\log(1-x).
\end{multline*}
Recall that $\mathsf g(x)=-\mathsf D\log x+n(U^{q_+}(x)-\kappa_+)$ for $x>0$. Add $\mathsf W_N$ to this and use \eqref{eq:qplus} for $0<x<1$ and \eqref{eq:gneg} for $-1<x<0$. The terms in $\log(1-x)$
cancel, the exponents of $1+x^2$ leave $2$, and the coefficient of $\log|x|$ is $75N-1-22N-\mathsf D=0$ in the first case and $75N-1-22N-53N=-1$ in the second. So
\begin{gather}
 \mathsf W_N(x)+n\mathsf R_*(x)+\mathsf g(x)\le nF_+(x)+2\log(1+x^2),\label{eq:rowpos}\\
 \mathsf W_N(x)+n\mathsf R_*(x)+\mathsf g(x)=nF_-(x)+2\log(1+x^2)-\log|x|.\label{eq:rowneg}
\end{gather}
Let $\mathsf R$ and $\mathsf C_0$ be upper bounds for the moved row and column fields at the given $\alpha$. The added fields are bounded, so they do not affect the following argument, which we write for $F_\pm$ and $\mathsf C$.
The term $-\frac12\log(1-x)$ of Lemma~\ref{lem:freetangent} and the term $-\log|x|$ of \eqref{eq:rowneg} are unbounded. They are controlled by the field itself. Let $H_0$ be the supremum of $F_+-\frac3{29}\log(1-x)-\frac{11}{29}\log|x|$ on $(0,1)$ and of $F_--\frac3{29}\log(1-x)-\frac{11}{29}\log|x|$ on $(-1,0)$. It is finite.
Indeed, after the subtraction $F_\sigma$ consists of constants, the term $-\frac{39}{29}\log(1+x^2)\le0$, the polynomial $\mathsf P$, the term $A_\sigma$ and
$U^{G_\sigma+R-2_\#C}(x+1/x)+m_\sigma\log|x|$. The term $A_\sigma$ is bounded because the auxiliary measures have bounded densities on circles. The measures $G_\sigma$, $R$ and $2_\#C$
are piecewise uniform with compact support, so their potentials are continuous on the real line. This settles the limits $x\to\pm1$, where $x+1/x\to\pm2$, and every compact subset of
$(-1,1)\setminus\{0\}$. As $x\to0$ we have $|x+1/x|\to\infty$. The measure $R-2_\#C$ has mass $0$ and $G_\sigma$ has mass $m_\sigma$, so $U^{G_\sigma+R-2_\#C}(x+1/x)+m_\sigma\log|x|$ tends to $0$.
Now split $nF=(n-\frac{29}6)F+\frac{29}6F$. Since $\frac{29}6\cdot\frac3{29}=\frac12$ and $\frac{29}6\cdot\frac{11}{29}=\frac{11}6\ge1$, and since $\log|x|\le0$,
\[
 nF-\tfrac12\log(1-x)-\log|x|\le(n-\tfrac{29}6)\mathsf R+\tfrac{29}6H_0 .
\]
The columns are treated in the same way. Let $H_1$ be the supremum of $\mathsf C(s)-\frac3{58}\log(1-s)$ on $(0,1)$. It is finite, because $U^C$ and $U^R(2s)$ are continuous and
$\log s\le0$. Since $\frac{58}3\cdot\frac3{58}=1$,
\[
 n\mathsf C(s)-\log(1-s)\le(n-\tfrac{58}3)\mathsf C_0+\tfrac{58}3H_1 .
\]
So every free point contributes at most $n\mathsf R+O(1)$ and every free column at most $n\mathsf C_0+O(1)$. 

Let now $y$ be a prescribed point. All prescribed points lie in $E_x$, so $y\ge98749/10^6$ or $y\le-19501/10^5$. Let $\mathsf H_+$, $\mathsf H_-$ and $\mathsf H_S$ be the fields of the prescribed rows in
the appendix. For $y>0$ outside $S$, \eqref{eq:qplus} and $53N=\mathsf D+1$ give
\[
 \mathsf w(y)+\mathsf g(y)\le \mathsf D\log2+n\mathsf H_+(y)+\log(1+y^2).
\]
For $y<0$, \eqref{eq:gneg} gives
\[
 \mathsf w(y)+\mathsf g(y)=\mathsf D\log2+n\mathsf H_-(y)-\log|y|+\log(1+y^2).
\]
For $y\in S$ the base point is $\xi=1/y$, so $\mathsf g=\mathsf D\log y+n(U^{q_+}(1/y)-\kappa_+)$. We have $\log|1/y-w|=-\log y+\log|1-yw|$, the measure $q_+$ has mass $1$, and $\int\log(1-yX)\,dG_+\le-\mathsf P(y)$
because all the moments are nonnegative. Hence $U^{q_+}(1/y)\le-\log y-\mathsf P(y)+\int\log|1-yw|\,d\nu_+(w)$. With $2\mathsf D=106N-2$ this gives
\[
 \mathsf w(y)+\mathsf g\le \mathsf D\log2+n\mathsf H_S(y)-2\log y+\log(1+y^2).
\]
The factor $e^{-N\hat\delta(\mathsf t(y))}$ of Proposition~\ref{prop:signed-goal} adds $\hat\delta(\mathsf t(y))/58$ to these fields. Let $\mathsf M$ be an upper bound for $\mathsf H_\pm+\hat\delta\circ\mathsf t/58$ on $E_W$ and for $\mathsf H_S+\hat\delta\circ\mathsf t/58$ on $E_W\cap(0,1)$. Since $\log(1+y^2)\le\log2$ and $|y|$ is bounded below on $E_x$, every prescribed row, with its factor $e^{-N\hat\delta}$, contributes at most
$\mathsf D\log2+n\mathsf M+O(1)$.

Now add \eqref{eq:scalar}, Lemma~\ref{lem:freetangent} multiplied by $n^2$, and half of Proposition~\ref{prop:gas}. With $n'=(1-\alpha)n$ and $\eta,\varepsilon\to0$, this gives
\[
 \log\bigl(\mathcal S\,|\det\Xi^0_S|\bigr)\le n^2\bigl(\mathsf B(\alpha)+o(1)\bigr),
\]
uniformly in the configuration, with
\begin{multline}\label{eq:budget}
 \mathsf B(\alpha)=\mathsf I-\tfrac32\alpha^2\log2+(1-\alpha)(\mathsf R+\mathsf C_0)+\alpha\mathsf M\\
 +\log2\Bigl(\tfrac{75}{58}(1-\alpha)+\tfrac12(1-\alpha)^2+\tfrac{53}{58}\alpha\Bigr).
\end{multline}
Here $\mathsf I=\mathsf I_0(\tau)+\mathsf I_{\mathrm{gas}}$, and $\frac12\mathsf K_N/n^2$ tends to
\[
 \mathsf I_{\mathrm{gas}}=\kappa_++\tfrac{53}{58}(\kappa_--c_M)-\tfrac12\bigl(\tfrac{53}{58}\bigr)^2\log2-\tfrac12\bigl(I(q_+)+I(q_-)\bigr)-\frac1{2\cdot58^2}\int_0^{53}\tau^*,
\]
with $\kappa_-$ and $c_M$ as above. The three terms in $\log2$ of \eqref{eq:budget} come from the exponent
$n'(75N-1)+\binom{n'}2$ and from the term $\mathsf D\log2$ of each prescribed row. Appendix~\ref{app:trials} gives the reference measures and the bounds
\begin{equation}\label{eq:caps}
\begin{gathered}
 \mathsf I_0+\mathsf I_{\mathrm{gas}}\le2.366997,\qquad\mathsf M\le-0.976,\\
 c^{\mathsf R}_0=-2.579740254,\quad c^{\mathsf R}_1=-2.574009165,\\
 c^{\mathsf C}_0=-2.946694099,\quad c^{\mathsf C}_1=-2.888907917,\quad c^{\mathsf C}_2=-2.85421372,
\end{gathered}
\end{equation}
with $\kappa_-=5463909/10^8$ and $\kappa_+=14762657/10^8$, and the moved values
\begin{gather*}
 L_0\le-7.5421576\cdot10^{-5},\quad L_1\le-0.04943292757,\\
 Q_{00}\le1.6623\cdot10^{-7},\quad Q_{01}\le9.5232\cdot10^{-7},\quad Q_{11}\le1.002863\cdot10^{-3}.
\end{gather*}
Let $V_S(\alpha)=\mathsf B(\alpha)+U_0-\frac12\log2-c_2\alpha$, with $\mathsf I$, $\mathsf R$ and $\mathsf C_0$ the moved values above and $\log2$
replaced by an upper bound. On each of the two pieces $[0,\alpha_e]$ and $[\alpha_e,\alpha_1]$ it is a polynomial in $\alpha$ with rational
coefficients. We have $V_S(0)=-1.81\ldots\cdot10^{-5}$, and the cover of \S\ref{ssec:cases} proves $V_S(\alpha)\le\sigma\alpha$ on both pieces. The
inequality is tight near $\alpha=0.0033$, where $V_S(\alpha)/\alpha$ reaches $0.33867$.

\begin{proof}[Proof of Proposition~\ref{prop:signed-goal}]
Fix a configuration with $m\le\alpha_1n$. By \eqref{eq:expandS} it suffices to bound each of the at most $2^m$ terms
$\mathcal S\,|\det\Xi_S|$, at the cost of the factor $5^m$. By Proposition~\ref{prop:alternative} each term either satisfies the
bound directly, or is at most $499^n\mathcal S\,|\det\Xi^0_S|$. In the second case the final inequality gives
$\exp\bigl(n^2(-U_0+\frac12\log2+c\alpha+o(1))-N\sum_j\hat\delta(t_{a_j})\bigr)$, since $499^n=e^{o(n^2)}$ and $c=c_2+\sigma$.
The factor $\prod_l(1+\zeta s_l)^m$ of Lemma~\ref{lem:binomvand} is contained in the column field, and the factors $\mathcal B_\zeta$ give
$\prod_j\mathcal B_\zeta(k_j)$.
\end{proof}
